\chapter{PCSEL 建模常见失败模式}
\label{ch:failure-modes}

\section*{学习目标}
\begin{itemize}
  \item 学会按结论层级、物理模型和数值实现三层顺序诊断 PCSEL 建模失败。
  \item 理解常见错误为什么会在 PCSEL 中尤其致命，而不是把它们当作普通求解器小毛病。
  \item 掌握在 Lumerical 与 COMSOL 工作流中最常见的误判路径和最优先的修复动作。
\end{itemize}

\section*{先修知识}
建议已掌握 \cref{ch:workflow,ch:robustness,ch:method-compare,ch:worked-example} 中关于工作流、鲁棒性分析、方法分工和示例执行链的内容。

\section*{本章路线图}
排错不是附录工作，而是建模能力本身的一部分。本章不按“软件功能表”写，而按“证据链在哪里断了”来写。我们先说明为什么 PCSEL 的错误往往先发生在结论层级而不是求解器按钮上，再按冷腔误判、边界条件误判、有效折射率误用、均匀增益误用、有限尺寸不足、鲁棒性分析缺席、热漂移忽略和单一求解器过度解读几个主类展开。随后再把这些错误映射到 Lumerical 与 COMSOL 中最常见的症状，给出优先排错顺序。

\section{把本章当作工作流和例题的反向读法}
\cref{ch:workflow} 的核心任务是告诉读者“正确顺序该怎么走”，\cref{ch:worked-example} 则把这条顺序真正走了一遍。本章应当被当作这两章的反向读法：当项目出现异常时，不是重新从头盲目扫参数，而是先问自己，问题究竟出在工作流的哪一段，以及若把自己的结果放回 \cref{tab:workflow-example-failure,tab:worked-map,tab:worked-summary}，最先对不上的那一栏是什么。这样做的好处是，排错会从“软件经验主义”转成“证据链审计”。

从实践角度看，一个成熟的项目通常会同时保存三张表：工作流章节的阶段映射表、本章例题的中间结果汇总表，以及本章的失败症状表。前两张表告诉你项目应该长成什么样，第三张表告诉你它偏离时会以什么症状表现出来。三者合起来，才能形成真正可执行的建模闭环。

\section{PCSEL 中最危险的失败，不是算不出来，而是算错了还以为自己算对了}
很多仿真项目的失败是显性的，例如求解器不收敛、网格内存爆掉、时域仿真提前终止。这类问题虽然麻烦，但至少容易察觉。PCSEL 建模中更危险的失败是隐性的：求解器给出了平滑曲线和漂亮场图，但结论的层级已经错了，或者近似的边界已被悄悄越过。这样的结果最容易进入汇报和论文，因为它“看起来像真的”。

这也是为什么本章首先强调诊断顺序。遇到问题时，先不要立刻怀疑软件，也不要立刻去调求解器参数，而应先问：我现在试图下的结论，是否超出了当前模型本身的适用层级。如果这个问题没过，后面的所有数值修补都只是为错误结论涂上一层更平滑的外观。

\begin{IntuitionBox}
PCSEL 建模里最昂贵的错误，通常不是少跑了几轮仿真，而是拿着一条层级错误的证据链往下越走越远。先纠正“我到底在证明什么”，通常比先纠正网格更值钱。
\end{IntuitionBox}

\section{第一类失败：把冷腔结果直接读成器件结论}
这是全书最反复强调、也是实际中最常见的失败模式。症状通常包括：用被动 $Q$ 直接宣布低阈值；用单胞反射谱最窄特征直接宣布器件最低阈值模式；用有限阵列的近场图直接宣布阈值电流或线宽优越。

根因很明确：冷腔结果只包含开放电磁问题，而器件结论至少还需要材料增益、注入分布和温度场。若这些变量还没有进入模型，所谓“阈值电流更低”“高温更稳定”“线宽更窄”都只是推测，不是结论。

修复顺序也应当明确：先把结论降级到 L1 或 L2，即“候选模更优”“光学阈值要求更低”“有限阵列远场更理想”；随后再决定是否值得进入电学和热学阶段。真正成熟的研究者，不会因为结论降级而觉得丢分，因为降级本身就是严谨的一部分。

\begin{RigorousBox}
若模型中尚不存在 $N(\mathbf{r})$、$T(\mathbf{r})$ 或与工作点有关的增益/吸收回写，就不应输出器件级阈值电流、工作线宽窗口或高温稳定性结论。这不是保守，而是对模型适用范围的诚实描述。
\end{RigorousBox}

\section{第二类失败：边界条件和几何定义在不同求解器里并不一致}
这类失败尤其容易出现在跨软件或跨阶段工作流中。比如，单胞 RCWA 使用的是严格周期边界，有限阵列 FDTD 却因为复制方式或阵列边界定义不当，使真实几何与单胞假设不一致；又或者 COMSOL 的单胞频域模型采用了某一相位周期条件，而后续有限器件模型却没有保持同样的坐标与材料定义。

症状通常表现为：不同求解器得到完全不同的频率排序；网格一加密结果反而跳来跳去；远场主瓣方向和偏振在两个模型中无法对上。遇到这种症状，最常见的误判是以为“某个求解器不可靠”。更常见的真实原因，其实是两个模型解的根本不是同一个物理问题。

修复这类问题时，最有效的做法不是在完整模型上继续试错，而是先退回到一个最小对照算例：只保留一层、一个孔、一个清楚的周期方向和一组共享参数，然后逐步确认边界、坐标、单位和材料定义完全一致，再把复杂层栈加回来。

\section{第三类失败：把有效折射率近似当成严格三维模型}
有效折射率近似是前期筛选的好工具，但它极易被误用成“已经足够描述器件”。症状通常是：二维模型预测的频率、偏振和远场都很好，但一进入三维验证就发生明显偏移；或者二维模型显示目标模有极佳重叠，三维模型却发现模式被上接触层或金属邻近区强烈拉偏。

根因在于有效折射率只压缩了纵向模的部分影响，并不能保留全部上下辐射不对称、层间损耗差异和强矢量效应。在 PCSEL 中，只要纵向剖面显著参与模式选择，这个近似就必须被当作代理模型，而不是最终结构模型。

修复方法也很直接：把二维模型的角色重新定义为“快速筛选器”，只让它决定参数方向，不让它独立承担最终结论；一旦候选结构进入关键排序比较，就必须回到三维矢量模型检查结果是否仍成立。

\section{第四类失败：把均匀增益默认当作真实电注入分布}
这是从光学走向器件时最危险的习惯性错误之一。很多初学者会在 RCWA、FDTD 或 Wave Optics 模型里直接给有源层加一个均匀增益，并据此默认自己已经“考虑了器件激射”。实际上，这最多只是小信号线性框架下的\textbf{均匀增益代理问题}。真实电注入器件中的增益分布取决于电流扩展、电流拥挤、漏电、串联电阻和温升，而这些都不在均匀增益假设里。

症状通常是：纯光学模型预测某个目标模最有利，但一旦把电流路径和热源加入，热点和高注入区却恰好偏向竞争模。此时若继续坚持均匀增益模型，就会把问题误判成“实验偏差太大”。真正的修复动作，是尽快引入 CHARGE 或 Semiconductor 一类电学模型，把均匀增益的假设显式替换成 $N(\mathbf{r})$ 诱导的空间分布。

\section{第五类失败：有限尺寸建模不足，导致把单胞排序误读成器件排序}
PCSEL 的很多关键特性都和有限孔阵有关，例如边缘泄漏、横向包络、主模尺寸和远场主瓣形成。若模型停留在单胞或极小阵列，就容易高估无限周期结论对真实器件的适用性。症状包括：单胞预测模式很理想，但实际有限阵列的远场副瓣很强；或者当阵列尺寸从 $21\times21$ 改到 $31\times31$ 时，目标模排序居然翻转。

修复这类问题时，最重要的不是立即追求更大阵列，而是先做尺寸收敛曲线。也就是说，至少把阵列尺寸作为一个显式自变量来扫，观察目标模与竞争模的排序、远场主瓣和边缘损耗是否收敛。若连这个趋势都不稳定，就没有资格把结果上升为器件级结论。

对执行层面而言，更稳妥的最小判据是：至少连续比较三个阵列尺寸，并同时跟踪三类量
\begin{equation}
\omega_\star(N_a),\qquad
Q_\star(N_a),\qquad
\mathcal{F}_{\mathrm{FF},\star}(N_a),
\label{eq:array_size_convergence}
\end{equation}
其中 $N_a$ 表示阵列周期数。若目标模相对首个竞争模的排序、主瓣方向或主瓣半高全宽在这三个尺寸之间仍翻转，则当前阵列尺寸就还不够大。\cref{eq:array_size_convergence} 的目的不是把“多少个周期才够”固定成某个神奇数字，而是把“有限尺寸足够”改写成必须展示的收敛趋势。

\begin{RigorousBox}
\textbf{有限阵列尺寸收敛判据。}

若用最小辐射衰减长度近似，阵列半径至少应满足 $R_{\mathrm{array}}\gtrsim 3L_{\mathrm{rad}}$，其中 $L_{\mathrm{rad}}\approx 1/(2\alpha_{\mathrm{rad}})$。这不是严格定理，但足以提醒读者：高 $Q$、低辐射损耗模往往需要远大于直觉的横向尺寸才能显露“接近无限周期”的行为。对亮模或中等 $Q$ 模，三维全波仍可能可行；对极高 $Q$ 的 quasi-BIC 候选模，所需阵列尺寸会很快超出直接三维 FDTD 的现实能力，此时应转向 RCWA 单胞加横向包络或边缘修正的混合策略。
\end{RigorousBox}

从执行可行性角度，也可以记一个粗窗口：几十微米量级的孔阵通常仍由边缘泄漏主导；进入 \SIrange{100}{300}{\micro\meter} 后，有限阵列全波开始有讨论价值；若再上到数百微米甚至更大并且目标模仍是高 $Q$ 候选，就应把“FDTD 是否仍然可行”当成工作流决策，而不是默认继续堆网格。

\section{第六类失败：忽略热漂移，默认冷腔频率就是工作频率}
热效应之所以在 PCSEL 中危险，不只是因为它会让阈值变高，而是因为它会主动改写模式竞争。折射率上升、增益峰位漂移、自由载流子吸收变化和横向温度梯度，都会把原本在冷腔中占优的模式重新排序。症状通常是：冷腔仿真与低功率实验能对上，但一到连续波或高占空比工作就出现远场偏移、偏振不稳或模式跳变。

这一类失败最常见的错误应对，是继续在冷腔里调几何参数，希望“把结构调到足够好”来抵消热漂移。更有效的应对策略是承认问题已经进入 L4 层级：必须把热源、温度场和光学回写显式纳入工作流。若还停留在冷腔里调参，通常只会越调越远。

\section{第七类失败：过度相信单一求解器或单一图像}
不同求解器各有优缺点。RCWA 强于周期单胞开放散射，FDTD 强于宽带有限阵列响应，COMSOL 的频域有限元强于复杂材料和多物理场联动，CHARGE / Semiconductor 强于漂移-扩散电学，HEAT / Heat Transfer 强于温度场求解。任何单一求解器都不可能完美覆盖所有层级问题。如果一个结论只在单一求解器里成立，而无法用第二种独立框架解释其趋势，那么它就还不够稳健。

这里要防止两个极端。一个极端是“某个软件最强，所以只信它”；另一个极端是“两个软件数值不一样，所以一定有一个错”。更成熟的做法是比较不变量和趋势，例如模式排序、主辐射方向、对参数扰动的单调性，而不是强迫两个求解器在每一个绝对数字上完全一致。

\begin{ExampleBox}
若 Lumerical 的 RCWA 与 COMSOL 的单胞频域模型都显示某个微扰会增强上方辐射并拉开两支偏振模的排序，那么即便两者给出的绝对 $Q$ 不完全相同，这个趋势仍然是有价值的。反过来，若两者对趋势都不一致，就应优先回查几何、边界和材料定义，而不是急着宣布哪个软件“更可信”。
\end{ExampleBox}

\section{这些失败在 Lumerical 与 COMSOL 中通常长什么样}
为了让排错更贴近实战，\cref{tab:failure-toolmap} 把常见症状映射到两条常见软件链中最典型的表现形式。注意，这个表的目的不是软件评测，而是帮助你更快定位“这更像是层级错误、物理模型错误，还是数值实现错误”。

\begin{table}[htbp]
  \centering
  \footnotesize
  \caption{常见失败症状在 Lumerical 与 COMSOL 工作流中的典型表现。}
  \label{tab:failure-toolmap}
  \setlength{\tabcolsep}{4pt}
  \begin{tabularx}{\textwidth}{@{}p{1.9cm}Y Y p{1.7cm}@{}}
    \toprule
    症状 & Lumerical 中常见表现 & COMSOL 中常见表现 & 优先动作 \\
    \midrule
    目标模频率 / $Q$ 对网格异常敏感 & FDTD 网格或仿真时间窗不足，PML 过近，或 RCWA 阶数不足 & 有限元网格不足、PML 设置不当，或本征频率研究捕捉到伪模 & 先做单因子收敛 \\
    周期单胞与有限阵列结果矛盾 & RCWA 与 FDTD 的阵列复制、边界或材料版本不一致 & 单胞周期模型与有限模型的几何或边界不一致 & 先做最小对照算例 \\
    引入增益后“阈值”过于乐观 & 把均匀增益直接当真实注入，未接 CHARGE / HEAT & 在 Wave Optics 中直接加均匀增益，未做 Semiconductor / Heat Transfer 回写 & 降级为光学阈值结论 \\
    高电流下远场和偏振突然恶化 & 未把温度场回写到折射率和增益代理 & 热场与光场未联动，仍在用冷态材料参数 & 进入热-电-光回写 \\
    电学求解看似收敛但结果不合理 & 接触边界、掺杂或电流归一化设置有误 & Semiconductor 接口的材料参数、边界或单位不一致 & 做量纲和边界审计 \\
    \bottomrule
  \end{tabularx}
\end{table}

若还想更具体一些，可以这样记忆：
\begin{itemize}
  \item 在 Lumerical 光学链里，最常见的危险信号是“结果对网格、PML 或仿真时间窗过分敏感”；
  \item 在 COMSOL 光学链里，最常见的危险信号是“本征频率或频域解看起来很平滑，但其实抓到的是错误边界条件下的模式”；
  \item 在两条电热链里，最常见的危险信号都是“结果数值稳定，但物理上明显不自洽”，例如热点位置与电流路径常识不符。
\end{itemize}

\section{推荐排错顺序：先层级，再物理，再数值}
当结果不可信时，最有效的排错顺序通常是三步。

第一步，检查结论层级。问自己：我现在试图声称的东西，是否超出了模型当前可支持的层级。

第二步，检查物理模型。问自己：边界条件、几何、材料、增益、注入和热源是否与研究问题匹配，是否存在被遗忘的近似。

第三步，最后才检查数值实现。包括网格、时间步、傅里叶阶数、PML、求解容差和 study 设置。

这个顺序看似反直觉，因为很多人更愿意先调数值参数。但对 PCSEL 来说，若结论层级本来就错了，继续调数值只会让错误结果更平滑，而不会让它更接近真实物理。

\begin{PitfallBox}
看到“两个求解器数值不同”就立刻开始加密网格，往往不是最高收益动作。若它们解的本来就不是同一个几何、同一组边界或同一层级问题，再细的网格也不会让它们自动一致。
\end{PitfallBox}

\section*{本章小结}
PCSEL 建模中最危险的失败，并不是显性的求解器报错，而是层级错误、近似越界和物理问题错配导致的“看起来合理的错误结论”。最常见的失败模式包括：把冷腔结果读成器件结论、边界条件和几何定义跨求解器不一致、误用有效折射率近似、把均匀增益当真实电注入、有限尺寸建模不足、忽略热漂移，以及过度依赖单一求解器。排错的正确顺序应是先检查结论层级，再检查物理模型，最后才检查数值实现。只要守住这个顺序，很多看似复杂的仿真问题都会变得可诊断、可修复。换句话说，\cref{ch:workflow} 是正向走法，\cref{ch:worked-example} 是标准做法，而本章就是它们的反向审计表。

\section*{练习题}
\begin{enumerate}
  \item \basicq 从你最近的一次仿真中选一个“结果看起来对，但你并不完全相信”的案例，按本章的三步顺序写出排错计划。

  \item \basicq 说明为什么把均匀增益模型直接用于电注入 PCSEL 的器件结论会造成系统性误判。

  \item \midq 设计一个最小对照算例，用来验证你在两个不同求解器中的边界条件和材料定义是否一致。

  \item \hardq 结合 \cref{tab:failure-toolmap}，写出一个同时适用于 Lumerical 与 COMSOL 的项目级排错清单。
\end{enumerate}

\section*{建议继续阅读}
\begin{itemize}
  \item 结合 \cref{ch:workflow,ch:sim-appendix} 与 \cref{tab:workflow-example-failure} 阅读，把本章中的失败模式直接转成你自己的项目模板和调试清单。 这条阅读的重点是把本章的输出顺着主线接到后续模块，而不是停成孤立知识点。

  \item 若准备开始真正的器件建模，建议把本章与 \cref{ch:worked-example,tab:worked-map,tab:worked-summary} 对照阅读：前者告诉你哪里最容易错，后者告诉你正确顺序该怎么走，以及每一步应该先产出什么证据。 这条阅读的重点是把本章的输出顺着主线接到后续模块，而不是停成孤立知识点。
\end{itemize}


